\begin{tikzpicture}[
    x=1pt,y=1pt,font=\small,>=stealth,
    block/.style={draw=#1!75!black, fill=#1!12, rounded corners=9pt, line width=0.9pt, align=center},
    flow/.style={->, line width=0.95pt, draw=gray!75}
]
\node[block=blue, minimum width=108pt, minimum height=38pt] (base) at (96,98) {\textbf{无防御基线}};
\node[block=orange, minimum width=108pt, minimum height=38pt] (single) at (258,98) {\textbf{单层防御对照}};
\node[block=teal, minimum width=108pt, minimum height=38pt] (full) at (420,98) {\textbf{三层协同防御}};
\node[block=red, minimum width=128pt, minimum height=48pt] (metrics) at (618,98) {\textbf{评估指标}\\ASR / UR / HR\\误报率 / 可用性损失};
\draw[flow] (base.east) -- (single.west);
\draw[flow] (single.east) -- (full.west);
\draw[flow] (full.east) -- (metrics.west);
\node[anchor=west, font=\bfseries\small] at (18,162) {防御效果评估框架图};
\node[anchor=west, font=\scriptsize, text=gray!70!black] at (18,30) {用于替代未完成的结果图组, 明确后续对照实验的指标与比较路径};
\end{tikzpicture}
